%%--------------------------------------------------------------------------
%%--------------------------------------------------------------------------
\subsection{Examen de ITT-SAT, 2 de julio de 2025}

Se quiere construir un sitio web, \textbf{YogurtLandia}, que permita a los usuarios explorar y gestionar una lista de tiendas de helado de yogur en una ciudad, así como interactuar con ellas. La funcionalidad básica del sitio es la siguiente:

\begin{enumerate}
    \item En el sitio no hay cuentas de usuario. Toda la funcionalidad está disponible para cualquiera que lo visite, salvo para añadir una nueva tienda o añadir una reseña y puntuación, para lo que se necesitará un nombre de usuario.
    \item La página principal del sitio tiene un formulario que permitirá a los visitantes escoger un nombre de usuario. Puede haber varios visitantes con el mismo nombre. Una vez elegido el nombre, este aparecerá en todas las páginas del sitio con el texto ``explorando como \emph{nombre} en YogurtLandia'', y debajo habrá un botón para \emph{cambiar usuario} (esto es, dejar de usar ese nombre).
    \item La página principal del sitio mostrará un listado con las diez tiendas de helado de yogur más populares de la ciudad. Este listado incluirá, por cada tienda: el nombre, su dirección, sus sabores destacados (ej. natural, fresa, chocolate), si tiene opciones veganas y la puntuación media recibida.
    \item Al seleccionar una tienda, el visitante verá la página específica de dicha tienda. En esa página aparecerá la misma información que en la página principal y, además, una galería de fotos de helados o del local subidas por otros usuarios (en orden cronológicamente inverso), así como el nombre de usuario que subió la foto, y todas las reseñas.
    \item Cada visitante puede subir fotos y añadir reseñas y puntuaciones (habiendo elegido nombre de usuario o no) utilizando un formulario en la página de la tienda. Tras subir la foto, añadir la reseña o envair puntuación, el visitante volverá a ver la página de la tienda que estaba viendo, pero ahora con la nueva información. La puntuación será de 1 a 5 estrellas.
    \item Toda página de YogurtLandia contiene una imagen de un pixel servido por el servicio \emph{sweetstats.com}.
\end{enumerate}

\newpage

Teniendo en cuenta los requisitos anteriores, se pide:

\begin{enumerate}
    \item Diseña un esquema REST para proporcionar el servicio descrito. Se habrán de especificar los nombres de recurso empleados, y cómo reaccionará la aplicación cuando reciba los métodos POST o GET sobre esas URLs (no se usarán los métodos PUT o DELETE). \textbf{Muy importante:} Coloca la información en una \textbf{tabla}, con los nombres de los recursos en una columna, los métodos en otra, la descripción de lo que realizará la aplicación al recibirlos en la tercera, y el contenido de los datos de la petición, si los hay en la cuarta (se consideran datos de la petición a los que vayan, normalmente como \emph{query string}, en el cuerpo de la petición HTTP o tras el nombre de recurso en la primera línea de la cabecera). Escribe también un fichero similar al fichero urls.py de Django (aunque no es importante que se respete la sintaxis mientras se entienda y la estructura sea similar a la de Django), que refleje el esquema REST anterior. (1 punto)

    \item Describe el modelo de datos que necesitará esta aplicación. Define las tablas necesarias y los campos necesarios para la funcionalidad descrita. Hazlo de forma lo más similar posible a lo que tendrías que escribir en el fichero models.py en Django (aunque no es importante que se respete la sintaxis mientras se entienda el modelo de datos que propones). (1 punto)

    \item Describe las interacciones HTTP que ocurrirán entre el navegador y cualquier servidor web en el siguiente escenario: El escenario comienza cuando un visitante (que accede por primera vez al sitio, esto es, pone en el navegador la URL de la página principal del sitio). A continuación, después de ver la página principal, elige un nombre de usuario, y luego añade una nueva tienda de helado de yogur. El escenario termina cuando el visitante ve la página de la tienda que acaba de crear. (1 punto)

    \item Describe las interacciones HTTP que ocurrirán entre el navegador y cualquier servidor web en el siguiente escenario: El escenario comienza cuando un visitante (que ha elegido el nombre ``pepe'') llega a la página de una tienda de helado de yogur en concreto. A continuación, después de ver la página, añade primero una imagen y luego una puntuación. El escenario termina cuando el visitante ve la página de la tienda con la nueva imagen y la nueva puntuación. (1 punto)

    \item Hay ciertas tiendas, que si visitas su página y le das una puntuación de 5 estrellas, te invitan a una consumición gratis si vuelves a visitar la página de la tienda entre una semana y un mes después. Uno de los requisitos técnicos es que la información necesaria para poder implementar esta funcionalidad no se guarde en el servidor. ¿Cómo se podría llevar a cabo? (1 punto)

\end{enumerate}

En todos los escenarios, ten en cuenta que tu respuesta debe considerar toda la funcionalidad que ofrece el servicio, y permitir que ésta pueda proporcionarse. Diseña la aplicación de forma que envíe cookies al navegador {\bf sólo cuando sea necesario}.

En las respuestas donde describas interacciones HTTP indica para cada una de ellas claramente y en este orden:
    \begin{itemize}
    \item La primera línea de la petición HTTP
    \item Si lo hay, el contenido de la petición
    \item La primera línea de la respuesta HTTP
    \item Si lo hay, el contenido de la respuesta
    \item Una brevísima explicación de para qué se usa la interacción
    \item Tanto en la petición como en la respuesta, las cabeceras con cookies, si es que fueran necesarias para la funcionalidad del escenario que se está describiendo (incluyendo el aspecto que han de tener esas cabeceras). Si la cabecera con cookie va o no dependiendo de algún factor ajeno a tu aplicación, explica cuando irá y cuándo no, y cuál es ese factor.
    \end{itemize}

Además, asegúrate de que describes las interacciones HTTP en el orden en que ocurrirían en el escenario.

